% Exemplo: referências que escrevem o tipo do elemento: \autoref (hyperref) e \cref (cleveref)
% Compilar: latexmk -pdf autoref-cleveref.tex      Testado com: TeX Live 2024 ou mais recente
\documentclass[a4paper]{article}
\usepackage[T1]{fontenc}
\usepackage{lmodern}
\usepackage[brazilian]{babel}
%%BEGIN:cleverefpreambulo
\usepackage[hidelinks]{hyperref}   % links e \autoref
\usepackage{cleveref}              % depois do hyperref
% nomes usados pelo \cref/\Cref (singular, plural)
\crefname{section}{seção}{seções}
\Crefname{section}{Seção}{Seções}
\crefname{table}{tabela}{tabelas}
\Crefname{table}{Tabela}{Tabelas}
%%END:cleverefpreambulo

\begin{document}

%%BEGIN:cleverefuso
\section{Introdução}\label{sec:intro}
Com \verb|\autoref|: veja a \autoref{sec:metodo}. Com \verb|\cref|: veja a \cref{sec:metodo};
no início da frase, \verb|\Cref|: \Cref{sec:metodo} descreve o método. Vários rótulos de uma
vez: \cref{sec:intro,sec:metodo}.

\section{Método}\label{sec:metodo}
%%END:cleverefuso

\end{document}
